% Chapter 4 — Implementation

\chapter{Implementation}
\label{Chapter4}

This chapter provides a detailed account of the implementation of each major system component. It describes the synthetic data generation pipeline, the multi-encoder NER training procedure, the QLoRA large language model fine-tuning, the ChromaDB vector index construction, the ensemble and hybrid pipeline architectures, and the Streamlit clinical workstation.

%----------------------------------------------------------------------------------------

\section{Synthetic Clinical Case Generation}

\subsection{Motivation and Dataset Gap}

No existing publicly available dataset provides patient narratives annotated with Kentian symptom dimensions and BIO token labels. The closest resources --- the i2b2/n2c2 NLP challenge corpora \cite{uzuner2011challenges} --- annotate standard medical concepts (medications, problems, tests) rather than the holistic symptom modalities required for homeopathic case-taking.

To address this gap, we constructed a large-scale synthetic corpus using the LLaMA 3 8B model via the Ollama REST API as a data generator. The generation approach is validated by the \textit{Annot-LLM} paradigm \cite{he2023annollm}, which demonstrates that LLM-generated annotations achieve competitive quality with human annotation for specialized NLP tasks when supervised by structured prompts.

\subsection{Generation Pipeline}

For each of the 4,933 MIND chapter rubrics, the generation pipeline produces 4 clinical variations per rubric, for a total of 19,732 cases. Each generation call uses a \textbf{few-shot prompt template} specifying:

\begin{itemize}
  \item The target rubric path (e.g., \texttt{MIND > ANXIETY > evening > alone, when})
  \item The required output JSON schema: \{\code{narrative}, \code{tokens}, \code{entities}\}
  \item Three clinical demonstration examples with corresponding BIO annotations
  \item Explicit instructions for language variation (formal/informal registers, age groups, symptom severity levels)
\end{itemize}

The \textbf{BIO Auto-Tagger} (\code{src/data/bio\_tagger.py}) performs post-generation validation:
\begin{enumerate}
  \item Tokenizes the narrative using the same tokenizer as the downstream NER encoder.
  \item Aligns entity character offsets to token boundaries using \textit{subword alignment repair}.
  \item Flags and repairs any offset drift (misalignment between character span and token span).
\end{enumerate}

\subsection{Dataset Statistics}

The final validated dataset contains:

\begin{table}[h]
\caption[Synthetic Dataset Statistics]{Kent-AI Synthetic Dataset Partition Statistics.}
\label{tab:dataset}
\centering
\begin{tabular}{lrrrr}
\toprule
\tabhead{Partition} & \tabhead{Cases} & \tabhead{Rubrics} & \tabhead{Entities} & \tabhead{Split \%} \\
\midrule
Train & 15,786 & 4,933 & 80,049 & 80\% \\
Validation & 1,973 & 1,973 & 9,850 & 10\% \\
Test & 1,973 & 1,973 & 9,955 & 10\% \\
\midrule
\textbf{Total} & \textbf{19,732} & \textbf{4,933} & \textbf{99,854} & 100\% \\
\bottomrule
\end{tabular}
\end{table}

Offset drift rate: \textbf{0 out of 99,854 entities} (100\% BIO alignment integrity). Mean case token length: 73.1 tokens (range: 34--131). Mean entities per case: 5.76.

\subsection{Pilot Validation Study}

A 100-case stratified pilot sample (25 rubrics $\times$ 4 variations) was manually audited against the following quality criteria:

\begin{table}[h]
\caption[Pilot Validation Results]{Pilot Study Quality Metrics (100 cases, 25 MIND rubrics).}
\label{tab:pilot}
\centering
\begin{tabular}{lrrr}
\toprule
\tabhead{Metric} & \tabhead{Result} & \tabhead{Target} & \tabhead{Status} \\
\midrule
JSON Schema Validity & 100\% (100/100) & 100\% & Pass \\
BIO Offset Accuracy & 100\% (576/576 spans) & $\geq$98\% & Pass \\
Token/BIO Tag Alignment & 100\% (0 mismatches) & 100\% & Pass \\
Total Entities & 576 (avg 5.76/case) & $\geq$4.0/case & Pass \\
Mean Pairwise Jaccard Overlap & 13.10\% & $<$45\% & Pass \\
Lexical Diversity Score & 86.90\% & $>$55\% & Pass \\
\bottomrule
\end{tabular}
\end{table}

%----------------------------------------------------------------------------------------

\section{Multi-Encoder NER Training}

\subsection{Training Configuration}

All four encoder models were fine-tuned using an identical training configuration to enable fair comparison:

\begin{itemize}
  \item \textbf{Task Formulation}: Token classification with 15 BIO output classes (B-/I- prefix for each of 7 dimensions, plus O for non-entity tokens).
  \item \textbf{Optimizer}: AdamW with weight decay 0.01.
  \item \textbf{Learning Rate}: $2 \times 10^{-5}$ with linear warmup over 10\% of steps.
  \item \textbf{Batch Size}: 16 (gradient accumulation $\times$ 2 for effective batch size 32).
  \item \textbf{Maximum Epochs}: 15 with early stopping (patience = 3) on validation F1.
  \item \textbf{Precision}: FP16 mixed-precision (AMP) on NVIDIA RTX 5000 Ada.
  \item \textbf{Hardware}: Single GPU, CUDA 12.4, PyTorch 2.6.0.
\end{itemize}

The \textbf{Training Loop} (\code{src/models/trainer.py}) implements the full training pipeline with SeqEval evaluation callbacks, TensorBoard logging, and automatic checkpoint saving at the best validation F1.

\subsection{Label Schema}

Table \ref{tab:labels} defines the 15 BIO token labels:

\begin{table}[h]
\caption[BIO Label Schema]{15-Class BIO Token Label Schema for Kent-AI NER.}
\label{tab:labels}
\centering
\begin{tabular}{lll}
\toprule
\tabhead{Label} & \tabhead{Dimension} & \tabhead{Clinical Meaning} \\
\midrule
B-MENT, I-MENT & Mental/Emotional & Anxiety, fears, grief, anger, restlessness \\
B-LOC, I-LOC & Location & Anatomical region, organ, side \\
B-SEN, I-SEN & Sensation & Pain quality, feeling, sensation type \\
B-MOD\_AGG, I-MOD\_AGG & Aggravation & What makes it worse \\
B-MOD\_AMEL, I-MOD\_AMEL & Amelioration & What makes it better \\
B-CONC, I-CONC & Concomitant & Accompanying symptoms \\
B-TEMP, I-TEMP & Temporal & Time, periodicity, duration \\
O & Outside & Non-entity tokens \\
\bottomrule
\end{tabular}
\end{table}

%----------------------------------------------------------------------------------------

\section{QLoRA Large Language Model Fine-Tuning}

\subsection{Instruction Dataset Construction}

The \code{src/models/llm\_dataset.py} module transforms the 15,786 training cases into instruction-formatted examples following the chat template of each target model. Each training instance follows the structure:

\begin{lstlisting}[style=python, caption={LLM Training Instance Format}]
{
  "system": "You are a clinical NER system for homeopathic case-taking...",
  "user": "Extract all symptom entities from: {narrative}",
  "assistant": '{"entities": [{"text": "...", "label": "MENT", ...}, ...]}'
}
\end{lstlisting}

\subsection{QLoRA Training Configuration}

\begin{itemize}
  \item \textbf{Quantization}: 4-bit NormalFloat (NF4) via BitsAndBytes, \code{double\_quant=True}.
  \item \textbf{LoRA Rank}: $r = 16$, $\alpha = 32$, targeting all attention projection matrices (\code{q\_proj, k\_proj, v\_proj, o\_proj, gate\_proj, up\_proj, down\_proj}).
  \item \textbf{Dropout}: 0.05 on LoRA layers.
  \item \textbf{Training}: SFTTrainer (HuggingFace TRL), 3 epochs, batch size 4, gradient accumulation 8 (effective 32).
  \item \textbf{Packing}: True (bin-packs short sequences for GPU utilization efficiency).
  \item \textbf{Max Sequence Length}: 2,048 tokens.
\end{itemize}

Final adapter checkpoint sizes: BioMistral = 81MB, Qwen = 78MB (versus 14GB+ for full 7B model weights).

%----------------------------------------------------------------------------------------

\section{ChromaDB Vector Index Construction}

The \textbf{build\_embeddings.py} script indexes all 74,513 rubrics in batches of 256:

\begin{enumerate}
  \item Each rubric is formatted as: \texttt{<section> > <chapter> > <subrubric> > <leaf>} (e.g., \texttt{MIND > ANXIETY > evening > alone, when}).
  \item The \code{all-MiniLM-L6-v2} model encodes each rubric string to a 384-dim L2-normalized vector.
  \item Vectors are batch-upserted to the ChromaDB HNSW collection with metadata: \{\code{rubric\_id}, \code{section\_id}, \code{grade}, \code{remedy\_count}\}.
\end{enumerate}

\textbf{Index Build Performance}: 74,513 rubrics indexed in 931.14 seconds at 80.0 rubrics/sec. Total collection size: 74,513 entries. Query latency: $<$18ms (p99) on 384-dim HNSW cosine index.

%----------------------------------------------------------------------------------------

\section{Ensemble and Hybrid Pipeline}

\subsection{4-Model Majority Voting Ensemble}

The ensemble in \code{src/models/ensemble.py} operates as follows:
\begin{enumerate}
  \item Collect BIO tag sequences from all four encoder models.
  \item Convert BIO sequences to named entity spans (text, start, end, label).
  \item Group overlapping spans using IoU threshold = 0.5.
  \item Assign the majority-vote label to each span cluster.
  \item Reject spans with $<$2 model agreement (minimum consensus threshold).
\end{enumerate}

\subsection{Two-Stage Hybrid Pipeline}

The hybrid pipeline in \code{src/models/hybrid\_pipeline.py} extends the ensemble:
\begin{enumerate}
  \item \textbf{Stage 1}: Ensemble mines candidate entity spans with any model agreement ($\geq$1 model).
  \item \textbf{Stage 2}: LLaMA 3.1 LLM arbitrator receives the full narrative plus candidate spans and produces a JSON-structured final extraction with label assignments and confidence flags.
  \item \textbf{Fallback}: If LLM is unavailable, the ensemble output (Stage 1) is used directly.
\end{enumerate}

%----------------------------------------------------------------------------------------

\section{Streamlit Clinical Workstation}

\subsection{Application Architecture}

The Streamlit application (\code{src/dashboard/app.py}) follows a navigation-router pattern:
\begin{itemize}
  \item \textbf{Top Navigation}: Horizontal menu tabs (Case Workspace, Case Library, Rubric Explorer, Materia Medica, Methodology, Settings).
  \item \textbf{Case Workspace}: 4-step stepper (Intake, Review, Analysis, Report) managed by a persistent \code{st.session\_state} stepper index.
  \item \textbf{Persistent Case Bar}: Top status bar showing case ID, patient name, draft status, and totality score. Always visible.
  \item \textbf{30\% Right Rail}: Case Snapshot with live dimension coverage indicators.
\end{itemize}

\subsection{PDF Export}

The \code{src/dashboard/pdf\_exporter.py} generates a multi-page ReportLab PDF with:
\begin{itemize}
  \item Cover page: Patient vitals, consultation metadata, clinical disclaimer.
  \item 7-Dimension Totality Grid: Colour-coded table of extracted entities by dimension.
  \item Totality Score Chart: Matplotlib bar chart of top-10 candidate remedies with grade-weighted scores.
  \item Candidate Remedy Comparison: Detailed table with rubric coverage, grade breakdown, and similarity scores.
  \item Two-pass \code{NumberedCanvas}: Running headers and ``Page X of Y'' footers across all pages.
\end{itemize}

\subsection{Testing Infrastructure}

The testing suite comprises 195 unit and integration tests across 22 test modules, covering:
\begin{itemize}
  \item Data layer: \code{test\_kent\_db.py}, \code{test\_case\_store.py}, \code{test\_bio\_tagger.py}.
  \item Model layer: \code{test\_symptom\_ner.py}, \code{test\_resolver.py}, \code{test\_ensemble.py}.
  \item Pipeline layer: \code{test\_pipeline.py}, \code{test\_ranker.py}, \code{test\_evaluation.py}.
  \item Dashboard layer: \code{test\_dashboard.py}, \code{test\_dialogue\_manager.py}.
\end{itemize}

All tests run without GPU or Ollama dependency via mock backends, enabling fast CI execution (full suite: $\sim$25 seconds).
